
\documentclass{article}
\usepackage{colortbl}
\usepackage{makecell}
\usepackage{multirow}
\usepackage{supertabular}

\begin{document}

\newcounter{utterance}

\centering \large Interaction Transcript for game `bbh', experiment `formal\_fallacies', episode 160 with qwen3.5{-}4b.
\vspace{24pt}

{ \footnotesize  \setcounter{utterance}{1}
\setlength{\tabcolsep}{0pt}
\begin{supertabular}{c@{$\;$}|p{.15\linewidth}@{}p{.15\linewidth}p{.15\linewidth}p{.15\linewidth}p{.15\linewidth}p{.15\linewidth}}
    \# & \multicolumn{2}{c}{Player} && \multicolumn{2}{c}{Game Master} \\
    \hline

    \theutterance \stepcounter{utterance}  
    & & & \multicolumn{4}{p{0.6\linewidth}}{
        \cellcolor[rgb]{0.9,0.9,0.9}{
            \makecell[{{p{\linewidth}}}]{
                \texttt{\tiny{[P1$\langle$GM]}}
                \texttt{Distinguish deductively valid arguments from formal fallacies.} \\
\\ 
\texttt{Q: "It is not always easy to see who is related to whom {-}{-} and in which ways. The following argument pertains to this question: To begin with, Lesley is a close friend of Fernando. Moreover, being a close friend of Fernando or a schoolmate of Lowell is sufficient for being a great{-}grandfather of Leroy. It follows that Lesley is a great{-}grandfather of Leroy."} \\
\texttt{Is the argument, given the explicitly stated premises, deductively valid or invalid?} \\
\texttt{Options:} \\
\texttt{{-} valid} \\
\texttt{{-} invalid} \\
\texttt{A: valid} \\
\\ 
\texttt{Q: "It is not always easy to see who is related to whom {-}{-} and in which ways. The following argument pertains to this question: Whoever is not a great{-}grandfather of Clyde is a stepbrother of Brian. Being an ancestor of Dana is sufficient for not being a great{-}grandfather of Clyde. We may conclude: Everyone who is an ancestor of Dana is a stepbrother of Brian, too."} \\
\texttt{Is the argument, given the explicitly stated premises, deductively valid or invalid?} \\
\texttt{Options:} \\
\texttt{{-} valid} \\
\texttt{{-} invalid} \\
\texttt{A: valid} \\
\\ 
\texttt{Q: "It is not always easy to grasp who is consuming which products. The following argument pertains to this question: Every infrequent user of Paul Mitchell shampoo is either a rare consumer of Nioxin shampoo or a loyal buyer of Caress soap, or both. No regular consumer of Lush soap is a rare consumer of Nioxin shampoo and, in the same time, a loyal buyer of Caress soap. It follows that whoever is an infrequent user of Paul Mitchell shampoo is not a regular consumer of Lush soap."} \\
\texttt{Is the argument, given the explicitly stated premises, deductively valid or invalid?} \\
\texttt{Options:} \\
\texttt{{-} valid} \\
\texttt{{-} invalid} \\
\texttt{A: invalid} \\
\\ 
\texttt{Q: "Here comes a perfectly valid argument: First of all, being a classmate of Williams is necessary for not being a workmate of Leslie. Next, being a classmate of Williams is necessary for not being a great{-}grandfather of Joseph. Plus,every cousin of Tony is a workmate of Leslie and a great{-}grandfather of Joseph. Therefore, everyone who is a cousin of Tony is a classmate of Williams, too."} \\
\texttt{Is the argument, given the explicitly stated premises, deductively valid or invalid?} \\
\texttt{Options:} \\
\texttt{{-} valid } \\
\texttt{{-} invalid} \\
\texttt{A:} \\
            }
        }
    }
     \\ \\

    \theutterance \stepcounter{utterance}  
    & \multicolumn{4}{p{0.6\linewidth}}{
        \cellcolor[rgb]{0.561,0.737,0.561}{
            \makecell[{{p{\linewidth}}}]{
                \texttt{\tiny{[P1$\rangle$GM]}}
                \texttt{To determine the validity of this argument, let's break down the premises and the conclusion using logical notation:} \\
\\ 
\texttt{**1. Symbolize the Premises:**} \\
\texttt{*   **Premise 1:** Being a classmate of Williams (\$C\$) is necessary for not being a workmate of Leslie (\$\textbackslash{}neg W\$).} \\
\texttt{    *   Logical form: \$\textbackslash{}neg W \textbackslash{}rightarrow C\$ (If you are not a workmate, you must be a classmate).} \\
\texttt{*   **Premise 2:** Being a classmate of Williams (\$C\$) is necessary for not being a great{-}grandfather of Joseph (\$\textbackslash{}neg G\$).} \\
\texttt{    *   Logical form: \$\textbackslash{}neg G \textbackslash{}rightarrow C\$ (If you are not a great{-}grandfather, you must be a classmate).} \\
\texttt{*   **Premise 3:** Every cousin of Tony (\$T\$) is a workmate of Leslie (\$W\$) and a great{-}grandfather of Joseph (\$G\$).} \\
\texttt{    *   Logical form: \$T \textbackslash{}rightarrow (W \textbackslash{}land G)\$ (If you are a cousin, you are a workmate AND a great{-}grandfather).} \\
\\ 
\texttt{**2. Analyze the Conclusion:**} \\
\texttt{*   **Conclusion:** Everyone who is a cousin of Tony (\$T\$) is a classmate of Williams (\$C\$).} \\
\texttt{    *   Logical form: \$T \textbackslash{}rightarrow C\$.} \\
\\ 
\texttt{**3. Evaluate the Logic:**} \\
\texttt{*   From Premise 3, we know that for any cousin of Tony (\$T\$), they are **both** a workmate (\$W\$) and a great{-}grandfather (\$G\$).} \\
\texttt{*   Therefore, for a cousin of Tony, the condition "not being a workmate" (\$\textbackslash{}neg W\$) is **false**, and the condition "not being a great{-}grandfather" (\$\textbackslash{}neg G\$) is **false**.} \\
\texttt{*   Now look at Premise 1: \$\textbackslash{}neg W \textbackslash{}rightarrow C\$. Since the antecedent (\$\textbackslash{}neg W\$) is false for a cousin of Tony, the implication tells us nothing about whether \$C\$ is true or false. (False \$\textbackslash{}rightarrow\$ True is valid; False \$\textbackslash{}rightarrow\$ False is valid; but we don't know the consequent).} \\
\texttt{*   Similarly, look at Premise 2: \$\textbackslash{}neg G \textbackslash{}rightarrow C\$. Since the antecedent (\$\textbackslash{}neg G\$) is false for a cousin of Tony, this premise also tells us nothing about \$C\$.} \\
\texttt{*   The premises establish that if someone is *not* a workmate or *not* a great{-}grandfather, they *must* be a classmate. However, the argument claims that because cousins *are* workmates and *are* great{-}grandfathers, they *must* be classmates. This is a logical error known as **denying the antecedent**. Just because the sufficient conditions for being a classmate (not being a workmate/great{-}grandfather) are not met, it does not mean the classmate status is impossible or guaranteed. Cousins could be classmates for other reasons, or they could be non{-}classmates entirely. The premises do not force the conclusion to be true.} \\
\\ 
\texttt{**Example to prove invalidity:**} \\
\texttt{*   Let "Cousin of Tony" = Person X.} \\
\texttt{*   Let Person X be a Workmate of Leslie (\$W\$) and a Great{-}grandfather of Joseph (\$G\$). (Satisfies Premise 3).} \\
\texttt{*   Let Person X **not** be a Classmate of Williams (\$\textbackslash{}neg C\$).} \\
\texttt{*   Check Premise 1: Is \$\textbackslash{}neg W \textbackslash{}rightarrow C\$ true? Since \$W\$ is true, \$\textbackslash{}neg W\$ is false. False \$\textbackslash{}rightarrow\$ False is a true statement. (Premise holds).} \\
\texttt{*   Check Premise 2: Is \$\textbackslash{}neg G \textbackslash{}rightarrow C\$ true? Since \$G\$ is true, \$\textbackslash{}neg G\$ is false. False \$\textbackslash{}rightarrow\$ False is a true statement. (Premise holds).} \\
\texttt{*   Check Conclusion: Is \$T \textbackslash{}rightarrow C\$ true? Here \$T\$ is true, but \$C\$ is false. So the conclusion is false.} \\
\texttt{*   Since we have a scenario where all premises are true but the conclusion is false, the argument is **invalid**.} \\
\\ 
\texttt{A: invalid} \\
            }
        }
    }
    & & \\ \\

    \theutterance \stepcounter{utterance}  
    & & & \multicolumn{2}{p{0.3\linewidth}}{
        \cellcolor[rgb]{0.9,0.9,0.9}{
            \makecell[{{p{\linewidth}}}]{
                \texttt{\tiny{[GM$|$GM]}}
                \texttt{Target: invalid} \\
            }
        }
    }
    & & \\ \\

    \theutterance \stepcounter{utterance}  
    & & & \multicolumn{2}{p{0.3\linewidth}}{
        \cellcolor[rgb]{0.9,0.9,0.9}{
            \makecell[{{p{\linewidth}}}]{
                \texttt{\tiny{[GM$|$GM]}}
                \texttt{game\_result = LOSE} \\
            }
        }
    }
    & & \\ \\

\end{supertabular}
}

\end{document}
